% Ch4 WS3 -- acid-base reactions and titration. Block 4.
\documentclass{apchem-worksheet}

\apcunitword{Chapter}
\apcunit{4}
\apcunittitle{Types of Chemical Reactions and Solution Stoichiometry}
\apcdoctitle{Worksheet 3 \textbullet\ Acids, Bases, and Titration}
\apcsubtitle{Zumdahl \S4.8 \textbullet\ decide strong vs weak BEFORE writing the equation}

\begin{document}
\wsheader[Every net ionic equation here depends on whether the acid is
strong or weak --- state that classification first, then write the
equation.]

\problem[]{Classify each as strong acid, weak acid, strong base, or weak
base:}
\begin{parts}
  \item \ce{HClO4}: \answerblank[3.6cm]{strong acid}
  \item \ce{HC2H3O2}: \answerblank[3.6cm]{weak acid}
  \item \ce{Ba(OH)2}: \answerblank[3.6cm]{strong base}
  \item \ce{NH3}: \answerblank[3.6cm]{weak base}
  \item \ce{HBr}: \answerblank[3.6cm]{strong acid}
  \item \ce{HF}: \answerblank[3.6cm]{weak acid}
\end{parts}

\problem[]{Write the net ionic equation for each neutralization:}
\begin{parts}
  \item \ce{HNO3(aq) + KOH(aq)}
        \answerblank[6.5cm]{\ce{H+ + OH^- -> H2O(l)}}
  \item \ce{HF(aq) + NaOH(aq)}
        \answerblank[6.5cm]{\ce{HF + OH^- -> F^- + H2O}}
  \item \ce{HCl(aq) + NH3(aq)}
        \answerblank[6.5cm]{\ce{H+ + NH3 -> NH4+}}
\end{parts}

\problem[]{Explain why (a) and (b) above have different net ionic equations
even though both are ``acid plus strong base.''
\answerlines[3]{\ce{HNO3} is a strong acid, fully ionized, so the reacting
species is free \ce{H+}. \ce{HF} is weak and remains almost entirely
molecular, so the whole \ce{HF} molecule transfers its proton directly to
hydroxide.}}

\problem[]{Give the Bronsted--Lowry definitions and explain why they are
more general than the Arrhenius definitions. Use \ce{NH3} as your example.
\answerlines[3]{Bronsted--Lowry: an acid donates a proton, a base accepts
one. This covers \ce{NH3}, which contains no \ce{OH-} yet acts as a base by
accepting \ce{H+} to form \ce{NH4+} --- something the Arrhenius definition
cannot explain.}}

\problem[]{Titration calculations.}
\begin{parts}
  \item \SI{25.0}{\milli\liter} of \ce{HCl} requires
        \SI{31.5}{\milli\liter} of \SI{0.200}{\Molar} \ce{NaOH}. Find
        $[\ce{HCl}]$. \workspace[2.2cm]
        \keyonly{\begin{apcanswer}$n(\ce{OH-}) = (0.0315)(0.200) =
        \SI{6.30e-3}{\mole}$; 1:1 $\to [\ce{HCl}] = 6.30\times10^{-3}/0.0250
        = \SI{0.252}{\Molar}$\end{apcanswer}}
  \item \SI{20.0}{\milli\liter} of \SI{0.150}{\Molar} \ce{H2SO4} is
        titrated with \SI{0.100}{\Molar} \ce{NaOH}. What volume of base is
        required? \workspace[2.4cm]
        \keyonly{\begin{apcanswer}$n(\ce{H2SO4}) = (0.0200)(0.150) =
        \SI{3.00e-3}{\mole}$; diprotic $\to n(\ce{OH-}) =
        \SI{6.00e-3}{\mole}$; $V = 6.00\times10^{-3}/0.100 =
        \SI{60.0}{\milli\liter}$\end{apcanswer}}
  \item \SI{0.500}{\gram} of an unknown monoprotic acid requires
        \SI{22.8}{\milli\liter} of \SI{0.250}{\Molar} \ce{NaOH}. Find the
        acid's molar mass. \workspace[2.4cm]
        \keyonly{\begin{apcanswer}$n = (0.0228)(0.250) = \SI{5.70e-3}{\mole}$;
        $M = 0.500/5.70\times10^{-3} =
        \SI{87.7}{\gram\per\mole}$\end{apcanswer}}
\end{parts}

\problem[]{Define the equivalence point and explain how an indicator signals
it. \answerlines[2]{The equivalence point is where the moles of titrant
added exactly satisfy the reaction stoichiometry of the analyte. An
indicator is chosen to change colour at that point, giving a visible
endpoint.}}

\problem[]{A student titrating \ce{H2SO4} with \ce{NaOH} uses a 1:1 mole
ratio and reports a concentration exactly half the true value. Explain the
error. \answerlines[2]{\ce{H2SO4} is diprotic --- each mole supplies two
moles of \ce{H+}, so it takes twice as much \ce{NaOH}. Using 1:1 attributes
all that base to twice as much acid as is actually present.}}

\begin{spiralstrip}{Chapter 4 \textbullet\ blocks 1--3}
  \item Net ionic for \ce{AgNO3(aq) + KCl(aq)}:
        \answerblank[5cm]{\ce{Ag+ + Cl^- -> AgCl(s)}}
  \item Soluble? \ce{BaSO4} \answerblank[2.4cm]{no (rule 4)}
  \item $[\ce{Na+}]$ in \SI{0.25}{\Molar} \ce{Na2CO3}:
        \answerblank[2.6cm]{\SI{0.50}{\Molar}}
  \item Strong electrolyte, weak, or non: glucose?
        \answerblank[3cm]{nonelectrolyte}
  \item Moles in \SI{50.0}{\milli\liter} of \SI{0.400}{\Molar}:
        \answerblank[2.8cm]{\SI{0.0200}{\mole}}
\end{spiralstrip}

\end{document}
